%%%------------------------------------
%%%     DIPLOMARBEIT: Hauptdatei
%%%     
%%%     Stefan Grünwald (2000 - 2007)
%%%     Michael Fladischer (2010)
%%%------------------------------------

\documentclass[%
    %DIV12,                                 % Satzspiegelparameter 12
    %BCOR8mm,                               % Bindekorrektur 8mm
    %headinclude,                           % Berücksichtigung der Kopfzeile
                                            % bei der Ermittlung des Satzspiegels
    12pt,                                   % Schriftgrösse 12pt
    a4paper,                                % DIN A4
    oneside,                                % einseitig
    draft=false,                            % Muss entfernt werden, bevor Arbeit abgegeben wird!
    bibliography=totoc,                     % Literaturverzeichniss automatisch im TOC einbinden
    listof=totoc,                           % Listen automatisch in TOC einbinden
    index=totoc]                            % Index automatisch in TOC einbinden
    {scrbook}                               % KOMA-Script Klasse 'scrbook' statt 'book'

% Packete
\usepackage{ngerman}
\usepackage[utf8]{inputenc}                 % Laden der UTF8-Eingabekodierung
                                            % ACHTUNG: Dateien müssen UTF8 codiert sein!
\usepackage[T1]{fontenc}                    % EC codierte Fonts
\usepackage{ae}                             % PDF encoding Info
\usepackage{palatino}                       % Brotschrift: Palationo
\usepackage{helvet}                         % Serifenlose Schrift: Helvetica
\usepackage[left]{eurosym}                  % EURO links
\usepackage{paralist}                       % Beeinflussung von Listen
%\usepackage{schuster}                      % \clubpenalty = 10000
                                            % \widowpenalty = 10000
                                            % \displaywidowpenalty = 10000
\usepackage{tabularx}                       % Austreiben auf eine bestimmte Breite
\usepackage{longtable}                      % Mehrseitige Tabellen
\usepackage{ltxtable}                       % 'longtable' und 'tabularx'
\usepackage{array}                          % Zeilenumbruch in 'p'-Spalten
\usepackage{booktabs}                       % \toprule, \midrule und \bottomrule
\usepackage{dcolumn}                        %
\usepackage{blindtext}                      % Erzeugt Blindtext welcher beim Prüfen des Layouts behilflich sein kann


% Lädt die lokalen Einstellungen
% --------------------- Angaben zur Arbeit ---------------------------------
\def\campusType{BACHELORARBEIT}
\def\campusTitle{Analyse zur Anwendung und Verarbeitung von signierten (X)HTML-Inhalten im WWW}
\def\campusAuthor{Michael Fladischer}
\def\campusPersKnz{0810319004}
\def\campusCity{Graz}
\def\campusDate{\today}
\def\campusSubject{}
\def\campusKeywords{PGP,x509}


% Um das Dokument sowohl für 'latex' als auch für 'pdflatex'
% kompatibel zu machen, wird das Package 'ifpdf' geladen
\usepackage{ifpdf}
% für Encapsulated-Postscript und PDFTeX Grafiken
\ifpdf
    \usepackage[pdftex]{graphicx}
    \graphicspath{{images/pdf/}}
    \pdfcompresslevel=9
\else
    \usepackage[dvips]{graphicx}
    \graphicspath{{images/dvi/}}
\fi

% Zeilenabstand
%\usepackage{setspace}
%\doublespace                               % doppelzeilig oder
%\onehalfspacing                            % eineinhalbzeilig

\usepackage{tikz}
\usetikzlibrary{mindmap,trees}

\usepackage{layout}
\usepackage[nooneline,hang]{caption}
\renewcommand{\captionfont}{\small\itshape}


\renewcommand{\floatpagefraction}{.8}       % ,float page`, wenn diese mindestens
                                            % 80% der Seite einnehmen
                                            % default: .5
\renewcommand{\textfraction}{.1}            % neben den Abbildungen, die ,h`ere, ,t`op und
                                            % ,b`ottom plaziert werden, noch min. 10% Text
                                            % default: .2
\renewcommand{\topfraction}{.9}             % max. Anteil der Seite, die f\"{u}r Abbildungen,
                                            % die ,t`op plaziert werden, verwendet wird
                                            % default: .7
\renewcommand{\bottomfraction}{.6}          % max. Anteil der Seite, die f\"{u}r Abbildungen,
                                            % die ,b`ottom plaziert werden, verwendet wird
                                            % default: .3


%%%------------------------------------
%%%---- DEFINITION INHALTSUEBERSICHT --
%%%------------------------------------
\makeatletter
\newcommand\shorttoc{\bgroup
%  \def\l@section##1##2{\relax}
  \def\l@subsection##1##2{\relax}
  \def\l@subsubsection##1##2{\relax}
  \section*{Inhalts\"{u}bersicht}
  \@input{\jobname.toc}
\egroup}
\makeatother

%%%------------------------------------
%%%---- SILBENTRENNUNG ----------------
%%%------------------------------------
% vim: set tw=160:

\hyphenation{Ge-schäfts-pro-zes-sen
             Ge-schäfts-pro-zess-mo-dell
             Ca-non-ic-aliz-a-tion-Meth-od
}



%\renewcommand{\el}[1]{}                     % extra label leeren
\usepackage{makeidx}
\makeindex

%%%------------------------------------
%%%---- 'pdf'-ERSTELLUNG --------------
%%%------------------------------------
\usepackage[a4paper,                        % DIN A4
%            dvips,                         % Erzeugung durch dvips: nicht mit 'pdflatex' verwenden
            plainpages=false,
            pdfpagelabels,
            pdftitle={\campusTitle},
            pdfauthor={\campusAuthor},
            pdfsubject={\campusType, \campusDate, \campusSubject},
            pdfkeywords={\campusKeywords},
            bookmarks=true,                 % Lesezeichen im 'pdf' analog zum Inhaltsverzeichnis,
                                            % (2x latex)
            bookmarksopen=true,             % Bookmarks anzeigen...
            bookmarksnumbered=true,         % ...und nummerieren
            hyperindex=true,                % Indexeinträge als Hyperlinks
            colorlinks,                     % Farbige Links
            linkcolor=black,
            urlcolor=black,
            citecolor=black,
            pdfborder=0,
            unicode=true
]{hyperref}

%%%------------------------------------
%%%---- GLOSSAR -----------------------
%%%------------------------------------
\usepackage[toc,acronym]{glossaries} 
\loadglsentries{./misc/glossary}
\makeglossaries

%%%---- EO USEPACKAGE -----------------
%%%------------------------------------

\newcommand{\uu}[1]{\underline{#1}}
\newcommand{\ii}[1]{\textit{#1}}

%%%------------------------------------
%%%---- LAYOUT ------------------------
%%%------------------------------------

% Numerierungstiefe der Kapitel
\setcounter{secnumdepth}{3}%

% Numerierungstiefe des Inhaltsverzeichnisses
\setcounter{tocdepth}{3}%

% Abstand vor und nach Gleitobjekten im Text [h]
%\setlength{\intextsep}{12pt plus 2pt minus 2pt}    % default: {12pt plus 2pt minus 2pt}

% Abstand zwischen Gleitobjekten oben oder unten und Text [tb]
\setlength{\textfloatsep}{28pt plus 2pt minus 4pt}  % default: {20pt plus 2pt minus 4pt}

%\setlength{\abovecaptionskip}{0pt}

%%%---- PARAGRAPH ---------------------
\makeatletter
\renewcommand{\paragraph}{%
  \@startsection {paragraph}{4}{\z@}%
  {-3.25ex plus-1ex minus-.2ex}{1sp}%
  {\raggedsection \normalfont \sectfont \nobreak \size@paragraph \nobreak } % was: \reset@font \normalsize\bf
} \makeatother

%%%---- SUBPARAGRAPH ------------------
\makeatletter
\renewcommand{\subparagraph}{%
  \@startsection {subparagraph}{5}{\parindent}%
  {3.25ex \@plus 1ex \@minus .2ex}{-1em}%
  {\raggedsection \normalfont \sectfont \nobreak \normalsize\bf \nobreak } % was: \size@subparagraph
} \makeatother

\clubpenalty = 10000
\widowpenalty = 10000
% \displaywidowpenalty = 10000

%\hyphenpenalty=9999                     % verhindert häufige Silbentrennung
%\lefthyphenmin=3                        % bei Silbentrennung bleiben min. 3 Zeichen
                                         % in der oberen Zeile stehen
%\finalhyphendemerits=5000%
%\doublehyphendemerits=10000%
%\sloppy                                 % sollte unbedingt vermieden werden ;-)

%%%------------------------------------
%%%---- SEITENDEKLARATION -------------
%%%------------------------------------
\setlength{\textheight}{660pt}%
\setlength{\voffset}{40pt}%
\setlength{\hoffset}{20pt}%
%\setlength{\footskip}{20pt}%

%%%------------------------------------
%%%---- KOPF-, FUSSZEILE --------------
%%%------------------------------------
\usepackage{fancyhdr}                   % Laden des Pakets zum Layout
%\pagestyle{fancy}                       % der Kopf- und Fusszeilen
\renewcommand{\chaptermark}[1]{\markboth{\thechapter.\ #1}{}}
%\renewcommand{\sectionmark}[1]{\markright{#1}{}}
%\renewcommand{\headrulewidth}{0.4pt}    % Linie Kopfzeile
%\renewcommand{\footrulewidth}{0.4pt}   % Linie Fusszeile
%\lhead{\sffamily\small\leftmark}        % LeftHeader --> Kapitelname
%\chead{}                                % CenteredHeader
%\rhead{\sffamily\small\thepage}         % RightHeader --> Seitennummer
%\lfoot{\sffamily\small}                 % LeftFooter --> Name
%\cfoot{}                                % CenteredFooter
%\rfoot{\sffamily\small}                 % RightFooter --> Titel
%%%------------------------------------

%%%------------------------------------
%%%---- BEGIN OF DOCUMENT -------------
%%%------------------------------------

\begin{document}

%%%------------------------------------
%%%---- VORSPANN ----------------------
%%%------------------------------------
\frontmatter                                % Titel, Vorwort, Inhaltsverzeichnis

%%%---- TITELSEITE --------------------
\begin{titlepage}
\begin{sf}
\begin{center}
\huge \campusType \\

\vspace{2cm}%

\campusTitle \\
\vspace{3cm}%

\begin{normalsize}

ausgeführt am

\begin{figure} [h!]
        \centerline{\includegraphics[width=0.3\textwidth]{campus02}}
\end{figure}

%\vspace{4.5cm}%

Studiengang\\ 
Informationstechnologien und IT-Marketing\\

\vspace{3.5cm}%

%\begin{tabularx}{\linewidth}{ll}
Von: \campusAuthor \\
Pers. Kennz. \campusPersKnz
%\end{tabularx}

\vspace{3.5cm}%
\end{normalsize}
\end{center}

\begin{tabularx}{0.9\linewidth}{Xc}
\campusCity, am \campusDate & ............................. \\
         & Unterschrift\\
\end{tabularx}

\end{sf}
\end{titlepage}

\newpage

% Seitennummerierung auf Römisch setzen und bei Seite 1 beginnen
\pagenumbering{roman}
\setcounter{page}{1}

%%%---- EO TITELSEITE -----------------

\setlength{\parindent}{0em}

%%%---- EHRENWORT ---------------------
% vim: set tw=160:

\chapter*{Ehrenwörtliche Erklärung}

Ich erkläre ehrenwörtlich, dass ich die vorliegende Arbeit selbstständig und ohne fremde Hilfe verfasst, andere als die angegebenen Quellen nicht benützt und
die benutzten Quellen wörtlich zitiert sowie inhaltlich entnommene Stellen als solche kenntlich gemacht habe.\\

\vspace{1 cm} \hfill
\begin{tabular*}{4.4 cm}{c}
....................................\\
\campusAuthor\\
\end{tabular*}


%%%---- DANK --------------------------
% vim: set tw=160:

\chapter*{Danksagung}

.....

\vspace{1.5cm}

\begin{center}
\raggedleft{\campusAuthor}
\end{center}

\vspace{1.5cm}

\campusCity, am \campusDate


%%%---- KURZFASSUNG -------------------
% vim: set tw=160:

\chapter*{Kurzfassung}
\blindtext%

% Hack damit zwei Chapter auf eine Seite passen
\begingroup
\let\clearpage\relax
\chapter*{Abstract}
\endgroup
\blindtext%


% Gleichheitsgrundsatz
% vim: set tw=160:

\chapter*{Gleichheitsgrundsatz}
Aus Gründen der Lesbarkeit wurde in dieser Arbeit darauf verzichtet, geschlechtsspezifische Formulierungen zu verwenden. Jedoch möchte der Author ausdrücklich
festhalten, dass die bei Personen verwendeten femininen oder maskulinen Formen für beide Geschlechter zu verstehen sind.


%%%---- INHALTSVERZEICHNIS ------------
\tableofcontents

%%%------------------------------------
%%%---- HAUPTTEIL ---------------------
%%%------------------------------------
\setlength{\parskip}{3ex plus0.5ex minus0.5ex}

%%%---- SEITENLAYOUT ------------------
%\newpage \layout

\mainmatter
% Seitennummerierung auf Arabisch setzen und bei Seite 1 beginnen
\pagenumbering{arabic}
\setcounter{page}{1}

% Laden des eigentlichen Inhalts
% vim: set tw=160:

%!TEX root =  ../diplomarbeit.tex
% vim: set spell spelllang=de_at,en_us
%%%%%%%%%%%%%%%%%%%%%%%%%%%%%%%%%%%%%%%
%%%     KAPITEL 1
%%%     EINLEITUNG
%%%%%%%%%%%%%%%%%%%%%%%%%%%%%%%%%%%%%%%
\chapter{Einleitung}
\index{Einleitung}%
\label{chp:einleitung}%


%%%------------------------------------
%%%---- AUSGANGSSITUATION -------------
%%%------------------------------------

\section{Ausgangssituation}
\index{Ausgangssituation}%
\label{sec:einl:ausgangssituation}%

Vertraulichkeit und Integrität von digitalen Inhalten gewinnt im in einem wachsenden globalen Netzwerk immer stärker an Bedeutung. Dies lässt sich anhand der Verbreitung von \gls{ssl}-fähigen Webservern und den darauf eingesetzten Zertifikaten bestimmen. \gls{ssl} dient jedoch nur der Vertraulichkeit Integrität der gesamten Daten und Authentizität der beiden Peers einer unidirektionalen Komunikation, wie sie z.B. bei HTTP zum Einsatz kommt. Aktive (umleiten der Verbindung) und passive Angriffe (Lesen der übertragenen Daten) wierden dadurch verhindert, eine Prüfung der Integrität und Authentizität einzelner, festlegbarer Bereiche des Inhalts ist damit jedoch nicht möglich. \\
In anderen Bereihen der Kommunikation wie Email und \gls{im} gehören Verfahren zum Nachweis der Authentizität des Absenders und der Integrität der Nachricht bereits zum Leitsungsumfang vieler Anwendungsprogramme. Die Verfahren welche dabei meist zum Einsatz kommen sind \gls{pgp} und \gls{smime}. Sie bieten die Möglichkeit, digitale Signaturen auf Basis von asymmetrischen Verschlüsselungsverfahren zu erstellen. Dabei wird von Verfasser der Nachricht diese mit seiner digitalen Signatur versehen. Der Empfänger ist nun in der Lage, die in der Nachricht signierten Teile zu überprüfen. \\
Mit dem Aufkommen des sog. Web 2.0 entstand auf vielen Webseiten für die Besucher die Möglichkeit, eigene Inhalte einzubinden oder bestehende Inhalte weiter zu verwenden. Der Betreiber einer Seite wurde dadurch vom alleinigen Bereitsteller zu einem Verwalter von Informationen. In Bezug auf Authentifität, Integrität und Verbindlichkeit des Inhalts waren Besucher somit meist auf optische Merkmale zu Separation von redaktionellem und Benutzer-generiertem Inhalt. \\
Digitale Signaturen könnten genau jenes Bedürfnis nach kryptografisch sicherer Trennung von Inhalten in einem \gls{xhtml}-Dokument erfüllen.

%%%------------------------------------
%%%---- AUFGABENSTELLUNG --------------
%%%------------------------------------
\section{Aufgabenstellung}
\index{Aufgabenstellung}%
\label{sec:einl:aufgabenstellung}%

Die primäre Aufgabenstellung dieser Arbeit besteht darin, zu bestimmen, welche Anforderungen an eine Integration von digitalen Signaturen in \gls{xhtml}-Dokumente gestellt wird. Dies soll aus zweierlei Perspektiven betrachtet werden:
\begin{itemize}
    \item Wie können digitale Signaturen in \gls{xhtml} integriert werden, ohne dass der entsprechende W3C-Standard\cite{xhtml:w3c} erweitert oder angepasst werden?
    \item Welche Informationen über eine oder mehrere digitale Signaturen sollen dem Benutzer in seinem Browser dargestellt werden, wenn er eine Webseite mit signierten Inhalten aufruft. Auf welche Art und Weise lässt sich diese Information benutzergerecht darstellen?
\end{itemize}
Weiters werden die technischen Aspekte der Integration einer Signaturprüfung in einem Browser erläutert:
\begin{itemize}
    \item Welche Implementierung eignet sich für die Integration in einen Browser?
    \item Wie kann die Verteilung der öffentlichen Schlüssel sowie die Referenzierung auf diese gelöst werden? Welche der bestehenden Technologien eigent sich dafür am besten?
\end{itemize}

%%%------------------------------------
%%%---- ZIELSETZUNG -------------------
%%%------------------------------------
\section{Zielsetzung der Arbeit}
\index{Ziele}%
\label{sec:einl:ziele}%
Diese Arbeit wird die Eignung der vorhandenen Signaturverfahren für \gls{xml}, Möglichkeiten zur Schlüsselverteilung und Präsentation der Signaturen im Browser behandeln.
Nicht behandelt wird die Erstellung eines neuen Signaturverfahrens sowie eines neuen Signaturmodells für \gls{xml} oder die Erweiterung des bestehenden \gls{xhtml}-Standards. Auch ist eine Implementierung eines Browser-Plugins nicht im Rahmen dieser Arbeit vorgesehen.

%%%------------------------------------
%%%---- VORGEHEN ----------------------
%%%------------------------------------
\section{Vorgehen und Methodik}
\index{Methodik}%
\label{sec:einl:methodik}%


%%%------------------------------------
%%%---- AUFBAU  -----------------------
%%%------------------------------------
\section{Aufbau der Arbeit}
\index{Aufbau der Arbeit}%
\label{sec:einl:aufbau}%

%%%%%%%%%%%%%%%%%%%%%%%%%%%%%%%%%%%%%%%
%%%     KAPITEL 2
%%%     Theoretische Grundlagen
%%%     2007/10/17
%%%%%%%%%%%%%%%%%%%%%%%%%%%%%%%%%%%%%%%
\chapter{Theoretische Grundlagen}
\index{Theoretische Grundlagen}%
\label{chp:TheoretischeGrundlagen}%

Prozess Beschreibung

%%%------------------------------------
%%%---- Requirements Engineering -------------
%%%------------------------------------

\section{Requirements Engineering}
\index{Requirements Engineering}%
\label{sec:thg:req}%


\subsection{Definition}
\index{Definition}
\label{sec:thg:req:definition}%

\subsection{Analysemodelle}
\index{Analysemodelle}
\label{sec:thg:req:modelle}%


%%%------------------------------------
%%%---- Prozessportfoliomanagement -------------
%%%------------------------------------

\section{Prozessportfoliomanagement}
\index{Prozessportfoliomanagement}%
\index{PPM}%
\label{sec:thg:ppm}%


\subsection{Definition}
\index{Definition}
\label{sec:thg:ppm:definition}%


\subsection{Elemente des Projektportfoliomanagement}
\index{PPM!Elemente}
\label{sec:thg:ppm:elemente}%

\subsection{Auswirkungen auf das Unternehmen}
\index{PPM!Unternehmen}
\label{sec:thg:ppm:auswirkungen}%

\subsection{Werkzeuge}
\index{PPM!Werkzeuge}
\label{sec:thg:ppm:tools}%


\subsection{Formulare, Analysetools}
\index{PPM!Formulare}
\index{PPM!Analysetools}
\label{sec:thg:ppm:forms}%

% vim: set tw=160:

\chapter{Signaturen mit XML-DSig}
\index{Signaturen mit XML-DSig}
\label{chap:SignaturenXMLDSig}

\section{Struktur einer Signatur}

\section{Referenzen}
\index{Referenzen}
\label{sec:Signaturbindung:Referenzen}

\section{Transformationen}
\index{Transformationen}
\label{sec:Signaturbindung:Transformationen}

\subsection{XPath}
\index{XPath}
\label{sec:Signaturbindung:Transformationen:XPath}

\subsection{XSLT}
\index{XSLT}
\label{sec:Signaturbindung:Transformationen:XSLT}

\section{Schlüsselinformation}

\section{Bindung an ein XHTML-Dokument}
Eine Signatur mit \gls{xml-dsig} hält ihre Verweise zu den signierten Inhalten über die Reference-Elemente. Im Fall der Verwendung innerhalb eines
\gls{webbrowser} muss das aufgerufene \gls{xhtml}-Dokument einen Verweis auf die Signatur beinhalten. Dies ist notwendig, da die Information über die Bereiche,
welche signiert wurden, nur in der Signatur selbst festgehalten werden. Ruft ein \gls{webbrowser} ein \gls{xhtml}-Dokument ab, so muss ihm bekannt gemacht
werden, dass für dieses Dokument Signaturinformationen vorhanden sind. Sind diese Informationen verfügbar gemacht worden, kann der \gls{webbrowser} daraus die
signierten Bereiche ermitteln.

Um solche Verweise auf Signaturen an die \gls{xhtml} binden zu können, bieten sich mehrere Verfahren an.

\subsection{HTTP-Header}
Eigener X-HTTP-Header wird vom Webserver ausgeliefert, worin über \glspl{url} auf Signatur-Dateien verwiesen wird. Keine Auswirkung auf XHTML, jedoch praktische Beschränkung
auf Anzahl der \glspl{url}?

\subsection{\xmlelem{link}-Element}
Der \gls{xhtml}-Standard definiert das \xmlelem{link}-Element, um Relationen zu anderen Dokumenten zu definieren \cite{xhtml:oreilly}. Es ist Teil des
\xmlelem{head}-Bereichs und kann dort mehrfach eingesetzt werden. Ein einzelnes \xmlelem{link}-Element definiert genau eine Relation. Bezüglich des Verhaltens
der \gls{webbrowser} beim Verarbeiten eines \xmlelem{link}-Elements macht der \gls{xhtml}-Standard keine verbindlichen Aussagen. Es wird aber darauf
hingewiesen, dass dieses Element speziell für zukünftige, noch nicht definierte Verhaltensmuster bei \gls{webbrowser} vorgesehen ist.

Es ist damit möglich, über das \xmlattr{href}-Attribut des \xmlelem{link}-Elements den Verweis auf die \gls{url} einer Signaturinformation zu realisieren. Da das
Element auch mehrfach in einem \gls{xhtml}-Dokument verwendet werden kann, lassen sich damit auch mehrere unterschiedliche Signaturinformationen einbinden.

\lstinputlisting[language=HTML,caption={Verwendung von \xmlelem{link}},label=lst:link,emph={link,href},emphstyle=\color{red}\bfseries\emph]{source/link.html}

\subsection{XML-Namespaces}
Direktes Einbetten der Signatur in das Dokument selbst.

\lstinputlisting[language=HTML,caption={Verwendung von XML-Namespaces},label=lst:xmlns,emph={Signature},emphstyle=\color{red}\bfseries\emph]{source/xmlns.html}

\subsection{XInclude}
XInclude ist eine neue Technologie, die vom W3C entwickelt wurde, um mehrere \gls{xml}-Dokumente in ein gemeinsames Dokument zusammenzuführen \cite{xml:oreilly}.
Dabei wird das \xmlelem{include}-Element aus dem reservieren XML-Namespace für \gls{xinclude} benutzt um über \gls{uri} externe \gls{xml}-Dokumente während der
Auswertung des Dokuments einzubinden.

\lstinputlisting[language=HTML,caption={Verwendung von XInclude},label=lst:xinclude,emph={include,xi},emphstyle=\color{red}\bfseries\emph]{source/xinclude.html}

% vim: set tw=160:

\chapter{Anwendbare Signaturverfahren}
\index{Signaturverfahren}
\label{chap:Signaturverfahren}

\section{Überblick}

\section{PGP}
\index{PGP}%
\label{sec:Signaturverfahren:PGP}%

\subsection{Verbreitung}

\subsection{Anwendung}

\subsection{Evaluierung}

\section{X.509}
\index{X.509}%
\label{sec:Signaturverfahren:x509}%

\subsection{Verbreitung}

\subsection{Anwendung}

\subsection{Evaluierung}

\section{Auswahl eines Verfahrens}


\chapter{Anforderungskatalog für ein Browser-Plugin}
\index{Anforderungskatalog}
\label{chp:Anforderungskatalog}

\section{Requirements Engineering}
\index{Requirements Engineering}
\label{sec:Anforderungskatalog:req}


\subsection{Definition}
\index{Definition}
\label{sec:Anforderungskatalog:req:definition}

\subsection{Analysemodelle}
\index{Analysemodelle}
\label{sec:Anforderungskatalog:req:modelle}%




%%%------------------------------------
%%%---- ANHANG ------------------------
%%%------------------------------------
\appendix

%%%------------------------------------
%%%---- NACHSPANN ---------------------
%%%------------------------------------
\backmatter %%% Bibliographie, Index

%%%------------------------------------
%%%     ABKUERZUNGSVERZEICHNIS
\clearpage%
\printglossaries
%%%------------------------------------
%%%     ABBILDUNGSVERZEICHNIS
\clearpage%
\listoffigures%
%%%------------------------------------
%%%     TABELLENVERZEICHNIS
\clearpage%
\listoftables%
%%%------------------------------------
%%%     LITERATURVERZEICHNIS
\clearpage%
\bibliographystyle{alphadin}%
\bibliography{./bibtex/literatur}%
%%%------------------------------------
%%%     INDEX
\clearpage%
\printindex
% vim: set tw=160:

\chapter*{}

\mbox{}

\foreigntextquote{german}[Michael Fladischer]{Have you tried turning it off and on again?}
%

%%%------------------------------------
%%%     EO DOCUMENT
%%%------------------------------------
\end{document}
